\section{第5章 定积分}
\subsection{定积分的概念与性质}
\begin{exercise}{557}{积分变量与参数求导}
设常数 $-1<a<b<1$，记 $I=\int_a^b\arcsin t\dd t$。分别求 $I$ 对无关变量 $x$、上限 $b$、下限 $a$ 的导数。若 $a=a(x),b=b(x)$ 为可导函数，写出 $\frac\dd{\dd x}\int_{a(x)}^{b(x)}\arcsin t\dd t$，说明哑变量命名不影响结果。
\end{exercise}
\begin{exercise}{558}{几何意义与尺度变换}
利用定积分的几何意义求 $\int_{-2}^{2}\sqrt{4-x^2}\dd x$，并推广计算 $\int_{-r}^{r}\sqrt{r^2-x^2}\dd x$（$r>0$）。若仅将积分区间改为 $[0,r]$，解释结果与原积分的关系。
\end{exercise}
\begin{exercise}{559}{黎曼和与端点影响}
利用定积分定义求 $\lim_{n\to\infty}\sum_{k=1}^{n-1}\frac{k}{n^2}$。对 $a>0$，再求 $\lim_{n\to\infty}\sum_{k=1}^{\lfloor an\rfloor}\frac{k}{n^2}$，说明取整及有限个端点项为何不影响极限。
\end{exercise}
\begin{exercise}{560}{黎曼和与积分夹逼}
将 $\sum_{k=1}^n\frac1{n+k}$ 写成黎曼和并求其极限。再利用 $f(x)=1/(1+x)$ 的单调性，证明该和与其极限的差的绝对值不超过 $1/(2n)$。
\end{exercise}
\begin{exercise}{561}{黎曼和与误差控制}
利用定积分定义求 $S_n=\sum_{k=1}^n\frac{n}{n^2+k^2}$ 的极限。将 $S_n$ 与相应积分比较，确定误差符号，并证明误差的绝对值不超过 $1/(2n)$。
\end{exercise}
\begin{exercise}{562}{积分比较与严格大小关系}
不计算原函数，分别比较以下各组积分：$\int_0^1x^2\dd x$ 与 $\int_0^1x^3\dd x$；$\int_0^1x\dd x$ 与 $\int_0^1\tan x\dd x$；$\int_1^2\ln x\dd x$ 与 $\int_1^2\ln^2x\dd x$；$\int_0^1e^x\dd x$ 与 $\int_0^1(1+x)\dd x$。给出每组严格不等式的依据，并指出端点相等是否影响结论。
\end{exercise}
\begin{exercise}{563}{根式积分的比较与区间影响}
不求原函数，比较 $I_1=\int_0^1\sqrt{1+x^2}\dd x$ 与 $I_2=\int_0^1\sqrt{1+x^3}\dd x$，证明严格大小关系。若积分区间同时改为 $[1,2]$，关系如何变化？说明理由。
\end{exercise}
\begin{exercise}{564}{积分估值与严格不等式}
不求原函数，先给出 $\int_{\pi/4}^{\pi}(1+\cos^2x)\dd x$ 的严格上、下界，所用界仅由被积函数最值与区间长度构成。再计算准确值，说明被积函数在有限个点达到最值为何不能使估值取等号。
\end{exercise}
\begin{exercise}{565}{带权积分估值}
不求原函数，证明 $\frac1{10\sqrt2}<\int_0^1\frac{x^9}{\sqrt{1+x}}\dd x<\frac1{10}$。把 $x^9$ 改为 $x^n$（$n$ 为非负整数），给出相应严格界，并求 $\lim_{n\to\infty}(n+1)\int_0^1\frac{x^n}{\sqrt{1+x}}\dd x$。
\end{exercise}
\begin{exercise}{566}{三角积分估值与单调性}
设 $I=\int_{\pi/4}^{\pi/2}\frac{\sin x}{x}\dd x$。证明 $\frac{\sqrt2}{4}\le I\le\frac{\sqrt2}{2}$，并判断两端能否取等号。再利用 $\sin x/x$ 的单调性改进下界。
\end{exercise}
\begin{exercise}{567}{积分中值定理与罗尔定理}
设 $a<b$，$f$ 在 $[a,b]$ 上连续、在 $(a,b)$ 可导，且 $\frac2{b-a}\int_a^{(a+b)/2}f(x)\dd x=f(b)$。证明存在 $\xi\in(a,b)$ 使 $f'(\xi)=0$，并说明可先在左半区间选出哪个中值点，再在哪个区间应用罗尔定理。
\end{exercise}
\begin{exercise}{568}{积分平均值的单调性}
设 $f$ 在 $[a,b]$ 上连续、在 $(a,b)$ 内可导，且 $f'\le0$。证明
\[
F(x)=\frac1{x-a}\int_a^xf(t)\dd t
\]
在 $(a,b)$ 内满足 $F'(x)\le0$。对固定的 $x_0\in(a,b)$，证明 $F'(x_0)=0$ 当且仅当 $f$ 在 $[a,x_0]$ 上为常数。
\end{exercise}
\begin{exercise}{569}{函数平均值与中值点}
求 $f(x)=\sqrt{4-x^2}$ 在 $[-2,2]$ 上的平均值，并求使函数值等于平均值的全部点。将半径 $2$ 改为 $r>0$ 后，求相应中值点与 $r$ 的比例关系。
\end{exercise}
\begin{exercise}{570}{积分定义与黎曼和}
设 $f$ 在 $[0,1]$ 上连续。用定积分表示极限
\[
\lim_{n\to\infty}\sum_{k=1}^n\frac{n}{n^2+k^2}f\!\left(\frac{k}{n}\right).
\]
当 $f(x)=x$ 时求极限值；当 $f(x)\ge0$ 且上述极限为零时，证明 $f\equiv0$。
\end{exercise}
\begin{exercise}{571}{积分不等式与等号条件}
设 $f$ 在 $[0,1]$ 上连续，且 $0\le f(x)\le1$、$\int_0^1f(x)\dd x=\frac12$。证明
\[
\frac14\le\int_0^1f^2(x)\dd x<\frac12,
\]
并完整讨论左端等号条件及右端不能取等号的原因。
\end{exercise}
\subsection{微积分基本公式}
\begin{exercise}{572}{变积分限与哑变量}
设 $F(a,b)=\int_a^b\sin t^2\dd t$。分别求其对无关变量 $x$、上限 $b$、下限 $a$ 的导数。再令 $a=x,b=x^2$，求复合函数的导数，说明不能把被积函数中的哑变量误作外部变量。
\end{exercise}
\begin{exercise}{573}{积分上限函数与凹凸性}
设 $y=\int_0^{x^2}\sqrt{1+t^2}\dd t$。求 $y'$、$y''$，判断 $y$ 的奇偶性及单调性，并判断是否存在拐点。
\end{exercise}
\begin{exercise}{574}{双变限积分的求导}
设 $y=\int_x^{x^2}\frac{\dd t}{\sqrt{1+t^2}}$。求 $y'$，并在 $x\ge0$ 上确定 $y$ 的零点与符号。计算 $y'(0)$、$y'(1)$，解释积分方向对符号判断的影响。
\end{exercise}
\begin{exercise}{575}{变限积分极限与高阶小量}
求 $\lim_{x\to0}\frac{\int_0^x\cos t^2\dd t}{x}$。进一步求 $\lim_{x\to0}\frac{x-\int_0^x\cos t^2\dd t}{x^5}$，并说明使用的局部展开及余项阶数。
\end{exercise}
\begin{exercise}{576}{变限积分极限与泰勒展开}
求 $\lim_{x\to0}\frac{\int_0^x\arctan t\dd t}{x^2}$，以及 $\lim_{x\to0}\frac{\int_0^x\arctan t\dd t-x^2/2}{x^4}$。论证从两侧趋于零时所得结果一致。
\end{exercise}
\begin{exercise}{577}{变限积分商的极限}
求
\[
\lim_{x\to0}\frac{\left(\int_0^xe^{t^2}\dd t\right)^2}{\int_0^xte^{2t^2}\dd t}.
\]
先证明分母在充分小的非零 $x$ 处不为零；再将分子中的指数 $t^2$ 改为 $at^2$、分母中的指数改为 $bt^2$，讨论实参数 $a,b$ 是否影响极限。
\end{exercise}
\begin{exercise}{578}{含参积分与连续性条件}
设 $f$ 在实轴上连续，求 $\lim_{x\to0}\frac{\int_0^xf(t)(x-t)\dd t}{x^2}$。分别用 $t=xu$ 与两次求导说明结果；再讨论将 $f$ 在原点连续弱化为在原点两侧有有限极限时，两侧极限相等的条件。
\end{exercise}
\begin{exercise}{579}{无穷远变限积分极限}
求 $\lim_{x\to+\infty}\frac{\int_0^x(\arctan t)^2\dd t}{\sqrt{1+x^2}}$。证明一般结论：若连续函数 $g(x)\to L$（$x\to+\infty$），则 $x^{-1}\int_0^xg(t)\dd t\to L$，并用该结论重新求本题极限。
\end{exercise}
\begin{exercise}{580}{隐函数与变限积分求导}
设 $\int_0^ye^t\dd t+\int_0^x\cos t\dd t=0$ 在原点附近确定隐函数 $y=y(x)$。求 $y'$、$y''$ 及其在 $x=0$ 的值，并确定包含 $0$ 的最大实定义区间。
\end{exercise}
\begin{exercise}{581}{高斯型变限积分的性质}
设 $f(x)=\int_0^xe^{-t^2/2}\dd t$，$x\in\mathbb R$。求其单调区间、奇偶性与凹凸区间；不用求出极限准确值，证明 $f$ 在实轴上有界。
\end{exercise}
\begin{exercise}{582}{含参变限积分与数据恢复}
设 $f$ 连续，$\varphi(x)=\int_0^{x^2}xf(t)\dd t$，且 $\varphi(1)=1$、$\varphi'(1)=5$。求 $f(1)$。若再知 $f$ 在 $1$ 处可导、$\varphi''(1)=9$，求 $f'(1)$，写清各次求导中的链式法则。
\end{exercise}
\begin{exercise}{583}{变限积分方程的唯一零点}
设 $f$ 在 $[a,b]$ 上连续且恒正，$a<b$。证明 $\int_a^xf(t)\dd t+\int_b^x\frac{\dd t}{f(t)}=0$ 在 $(a,b)$ 恰有一根。若 $f\equiv c>0$，求该根，并判断其关于 $c$ 的单调性。
\end{exercise}
\begin{exercise}{584}{基本公式与参数积分}
计算 $\int_{-1}^{1/2}e^{2x}\dd x$。对实参数 $a$，计算 $J(a)=\int_{-1}^{1/2}e^{ax}\dd x$，单独处理 $a=0$，并证明所得表达式在 $a=0$ 连续。
\end{exercise}
\begin{exercise}{585}{原函数与参数上限}
计算 $\int_0^2\frac{x}{\sqrt{1+x^2}}\dd x$。令 $G(a)=\int_0^a\frac{x}{\sqrt{1+x^2}}\dd x$，$a\in\mathbb R$，求 $G$ 并判断其奇偶性及最小值。
\end{exercise}
\begin{exercise}{586}{部分分式与积分取值}
计算 $\int_0^1\frac{\dd x}{x^2+5x+6}$。再令上限为 $a\ge0$，求积分关于 $a$ 的表达式，并确定其值域。
\end{exercise}
\begin{exercise}{587}{绝对值积分与参数最值}
计算 $\int_{-1}^3|2-x|\dd x$。对实数 $a$ 求 $H(a)=\int_{-1}^3|a-x|\dd x$ 的分段表达式，并求 $H$ 的最小值及取到最小值的全部 $a$。
\end{exercise}
\begin{exercise}{588}{绝对值积分与周期性}
计算 $\int_0^{2\pi}|\sin x|\dd x$。证明对任意实数 $a$ 与正整数 $n$，$\int_a^{a+n\pi}|\sin x|\dd x$ 与 $a$ 无关，并求其值。
\end{exercise}
\begin{exercise}{589}{分段函数的变限积分与可导性}
设 $f(x)=2x$（$0\le x\le1$），$f(x)=2+x$（$1<x\le2$）。求 $P(x)=\int_0^xf(t)\dd t$ 在 $[0,2]$ 上的表达式，判断 $P$ 在 $x=1$ 处的连续性和可导性，并求该处左右导数。
\end{exercise}
\begin{exercise}{590}{积分方程与自洽条件}
连续函数 $f$ 在 $[0,1]$ 上满足 $f(x)=\frac1{1+x^2}+\sqrt{1-x^2}\int_0^1f(t)\dd t$。求 $\int_0^1f(t)\dd t$ 及 $f$。若根式项前再乘实参数 $\lambda$，讨论连续解的存在性与唯一性。
\end{exercise}
\begin{exercise}{591}{积分方程与参数可解性}
设 $f$ 在 $[0,1]$ 上连续，满足 $f(x)=x\int_0^1f(t)\dd t-1$。求 $\int_0^1f(t)\dd t$ 及 $f$。若 $x$ 前的系数改为实数 $\lambda$，讨论方程何时存在唯一连续解。
\end{exercise}
\begin{exercise}{592}{复合函数恢复与换元}
设 $f$ 在 $[2,+\infty)$ 上满足 $f(x+1/x)=\frac{x+x^3}{1+x^4}$（$x>0$）。证明此条件唯一确定 $f$，求 $f$ 的表达式及 $\int_2^{2\sqrt2}f(u)\dd u$。说明 $x$ 与 $1/x$ 给出的函数值为何一致。
\end{exercise}
\begin{exercise}{593}{变限积分与极值}
设 $F(x)=\int_0^x(t^2-1)e^{-t}\dd t$，$x\ge0$。求 $F$ 的单调区间与极值，并计算 $\lim_{x\to\infty}F(x)$；由此确定 $F$ 在 $[0,+\infty)$ 上的零点个数。
\end{exercise}
\begin{exercise}{594}{积分方程与条件恢复}
设连续函数 $f$ 满足
\[
f(x)=1+\int_0^x(x-t)f(t)\dd t\qquad(x\in\mathbb R).
\]
证明 $f$ 二阶可导，求 $f$；再直接代回积分等式核验。
\end{exercise}
\subsection{定积分的换元积分法和分部积分法}
\begin{exercise}{595}{根式换元与定义域}
计算 $\int_4^9\frac{\sqrt x}{\sqrt x-1}\dd x$。对 $1<a<b$，求同一被积函数在 $[a,b]$ 上的积分，并说明当下限趋于 $1^+$ 时积分的变化。
\end{exercise}
\begin{exercise}{596}{根式换元与积分函数}
计算 $\int_0^1x\sqrt{3-2x}\dd x$。对 $0\le a\le3/2$，求 $J(a)=\int_0^ax\sqrt{3-2x}\dd x$，并确定 $J$ 在该区间上的最大值。
\end{exercise}
\begin{exercise}{597}{对数换元与定义域}
计算 $\int_1^e\frac{\dd x}{x\sqrt{1+\ln x}}$。将上限改为 $a>e^{-1}$，求积分表达式；再判断上限固定为 $1$、下限趋于 $e^{-1}$ 时的积分是否有限。
\end{exercise}
\begin{exercise}{598}{复合换元与单调条件}
计算 $\int_{-1}^0\frac{(1+x)e^x}{\sqrt{1+xe^x}}\dd x$。说明换元 $u=1+xe^x$ 在积分区间上是否合法；再求将下限改为 $a\in[-1,0]$ 后的积分及其关于 $a$ 的单调性。
\end{exercise}
\begin{exercise}{599}{三角换元与几何意义}
计算 $\int_0^{\sqrt2/3}\sqrt{2-9x^2}\dd x$，并用椭圆面积核验。进一步计算 $\int_0^{\sqrt2/3}x^2\sqrt{2-9x^2}\dd x$。
\end{exercise}
\begin{exercise}{600}{三角换元与无穷上限}
计算 $\int_1^{\sqrt3}\frac{\dd x}{x\sqrt{1+x^2}}$。求 $J(a)=\int_1^a\frac{\dd x}{x\sqrt{1+x^2}}$（$a\ge1$），并求 $\lim_{a\to+\infty}J(a)$。
\end{exercise}
\begin{exercise}{601}{根式积分与反三角函数}
计算 $\int_1^2\frac{\sqrt{x^2-1}}x\dd x$。将上限改为 $a\ge1$，求积分表达式；由此求 $\lim_{a\to+\infty}\left(\int_1^a\frac{\sqrt{x^2-1}}x\dd x-a\right)$。
\end{exercise}
\begin{exercise}{602}{根式积分与对称换元}
计算 $\int_{-1}^1\frac{\dd x}{1+\sqrt{1-x^2}}$。对 $-1\le a\le1$ 求 $J(a)=\int_{-a}^a\frac{\dd x}{1+\sqrt{1-x^2}}$，正确处理 $a<0$ 时的积分方向，并判断 $J$ 的奇偶性。
\end{exercise}
\begin{exercise}{603}{配方换元与圆面积}
计算 $\int_0^1\sqrt{2x-x^2}\dd x$。利用被积函数关于 $x=1$ 的对称性求 $\int_0^2x\sqrt{2x-x^2}\dd x$，并说明不必另求原函数的理由。
\end{exercise}
\begin{exercise}{604}{指数根式换元与积分方向}
计算 $\int_0^{-\ln2}\sqrt{1-e^{2x}}\dd x$。设 $J(a)=\int_{-a}^0\sqrt{1-e^{2x}}\dd x$（$a>0$），求 $J(a)$ 及 $\lim_{a\to+\infty}(J(a)-a)$，先由积分方向判断原积分的符号。
\end{exercise}
\begin{exercise}{605}{平移换元与分段积分}
设 $f(x)=1+x^2$（$x<0$），$f(x)=e^{-x}$（$x\ge0$）。求 $\int_1^3f(x-2)\dd x$；令 $J(c)=\int_1^3f(x-c)\dd x$，证明 $J$ 在 $c=2$ 可导并求 $J'(2)$。
\end{exercise}
\begin{exercise}{606}{分段函数的平移与奇偶消去}
设 $f(x)=xe^{x^2}$（$-1/2\le x\le1/2$），$f(x)=-1$（$x>1/2$）。求 $\int_{1/2}^2f(x-1)\dd x$。说明换元后如何利用奇偶性简化；若仅改变 $f(1/2)$ 的值，积分是否改变？给出理由。
\end{exercise}
\begin{exercise}{607}{卷积换元与函数恢复}
设连续函数 $f$ 满足 $\int_0^xf(x-t)\dd t=e^{-2x}-1$。求 $\int_0^1f(x)\dd x$ 及 $f(x)$，并证明此连续解唯一。
\end{exercise}
\begin{exercise}{608}{含参积分与二次求导}
设 $f$ 在实轴上连续，且 $\int_0^xf(x-t)t\dd t=e^x-x-1$。通过换元和微积分基本定理求 $f$，并直接核验原等式。说明推导时为何不需要预先假设 $f$ 可导。
\end{exercise}
\begin{exercise}{609}{奇函数积分与参数平移}
计算 $\int_{-\pi/2}^{\pi/2}\frac{x}{1+\cos x}\dd x$。将分子改为 $x+c$，求积分关于实参数 $c$ 的表达式，并确定积分为零的充要条件。
\end{exercise}
\begin{exercise}{610}{对数函数的奇偶性}
计算 $\int_{-\pi/2}^{\pi/2}\sin^2x\ln(x+\sqrt{1+x^2})\dd x$。先证明对数因子是奇函数；若将 $\sin^2x$ 改为任意连续函数 $g(x)$，说明积分只与 $g$ 的奇部还是偶部有关并证明。
\end{exercise}
\begin{exercise}{611}{对称换元与参数条件}
计算 $\int_{-1}^1\frac1{\sqrt{4-x^2}}\left(\frac1{1+e^x}-\frac12\right)\dd x$。若 $\frac12$ 改为实参数 $c$，求积分关于 $c$ 的表达式及其为零的充要条件。
\end{exercise}
\begin{exercise}{612}{反三角函数的对称关系}
计算 $\int_{-1}^1\cos x\arccos x\dd x$。证明对任意连续偶函数 $g$，$\int_{-1}^1g(x)\arccos x\dd x=\frac\pi2\int_{-1}^1g(x)\dd x$，并说明本题如何应用此式。
\end{exercise}
\begin{exercise}{613}{奇偶消去与积分排序}
设三积分的区间均为 $[-\pi/4,\pi/4]$，被积函数分别为 $\frac{\tan x}{1+x^4}+x^8$、$\sin^8x+\ln(x+\sqrt{1+x^2})$、$\tan^8x+e^x\cos x-e^{-x}\cos x$，对应积分为 $M,N,P$。不用求原函数，确定三者严格大小次序；证明奇函数部分消去后所用的点态不等式，并说明严格性的来源。
\end{exercise}
\begin{exercise}{614}{周期函数与绝对值}
对正整数 $n$，计算 $\int_0^{n\pi}\sqrt{1-\sin^2x}\dd x$。再求 $\int_a^{a+n\pi}\sqrt{1-\sin^2x}\dd x$，证明结果与实数 $a$ 无关，并指出去掉根号而直接写成 $\cos x$ 的错误。
\end{exercise}
\begin{exercise}{615}{周期表达式的瑕点与截断}
考察 $I=\int_0^\pi\sin^2x(\tan x+1)\dd x$。指出区间内部的瑕点，按反常积分定义判断 $I$ 是否存在；再求对称截断和 $\lim_{\varepsilon\to0^+}(\int_0^{\pi/2-\varepsilon}+\int_{\pi/2+\varepsilon}^{\pi})\sin^2x(\tan x+1)\dd x$，说明两种概念的区别。
\end{exercise}
\begin{exercise}{616}{对称换元恒等式与指数积分}
设 $f$ 在 $[0,1]$ 上连续。证明 $\int_0^{\pi/2}f(\sin x)\dd x=\int_0^{\pi/2}f(\cos x)\dd x$ 及 $\int_0^\pi xf(\sin x)\dd x=\pi\int_0^{\pi/2}f(\sin x)\dd x$。利用这两式计算 $I_n=\int_0^\pi\frac{x\sin^{2n}x}{\sin^{2n}x+\cos^{2n}x}\dd x$（$n$ 为正整数），证明结果与 $n$ 无关。
\end{exercise}
\begin{exercise}{617}{分部积分与参数恒等式}
计算 $\int_0^1\ln(1+x^2)\dd x$。对 $a>0$ 计算 $\int_0^a\ln(1+x^2)\dd x$，再对所得结果求导核验。
\end{exercise}
\begin{exercise}{618}{反三角函数的分部积分}
计算 $\int_0^{1/2}\arcsin x\dd x$。将上限改为 $a\in[-1,1]$，求积分表达式并判断其奇偶性；由所得结果求该积分在允许参数范围内的最大值。
\end{exercise}
\begin{exercise}{619}{分部积分与奇偶分解}
计算 $\int_{-1}^1x\arccos x\dd x$。分别用分部积分及 $\arccos(-x)=\pi-\arccos x$ 将积分化到 $[0,1]$ 上，说明两种计算的联系。
\end{exercise}
\begin{exercise}{620}{分部积分与有理积分}
计算 $\int_0^1\frac{\ln(1+x)}{(2-x)^2}\dd x$。先选择能使剩余积分化为有理函数的分部方式；再利用被积函数的符号核验所得结果为正。
\end{exercise}
\begin{exercise}{621}{循环分部积分}
计算 $\int_0^{\pi/2}e^{-x}\sin2x\dd x$。求 $F(a)=\int_0^ae^{-x}\sin2x\dd x$（$a\ge0$），并由 $F'(a)$ 判断 $a=\pi/2$ 是极大点还是极小点。
\end{exercise}
\begin{exercise}{622}{换元与分部积分的结合}
计算 $\int_0^1e^{-\sqrt x}\dd x$。对 $a>0$ 求 $\int_0^{a^2}e^{-\sqrt x}\dd x$，并求其在 $a\to+\infty$ 时的极限。
\end{exercise}
\begin{exercise}{623}{三角换元与端点处理}
计算 $\int_0^1\frac{x^2\arcsin x}{\sqrt{1-x^2}}\dd x$，在计算中说明 $x=1$ 处无界因子为何不导致发散；再求 $\int_0^1\frac{\arcsin x}{\sqrt{1-x^2}}\dd x$ 作为对照。
\end{exercise}
\begin{exercise}{624}{嵌套根式换元与分部积分}
计算 $\int_1^{16}\arctan\sqrt{\sqrt x-1}\dd x$。先说明 $u=\sqrt{\sqrt x-1}$ 将区间变为何区间，再用分部积分化为多项式积分；用被积函数取值范围给出结果应满足的粗略界。
\end{exercise}
\begin{exercise}{625}{倒数换元与指数分部积分}
计算 $\int_1^2x^{-3}e^{1/x}\dd x$。将上限改为 $a\ge1$ 求积分表达式，并判断该积分在 $a\to+\infty$ 时是否有有限极限，给出极限值。
\end{exercise}
\begin{exercise}{626}{根式换元与振荡符号}
计算 $\int_0^{\pi^2}\sqrt x\cos\sqrt x\dd x$。先换元化为多项式乘三角函数，再解释原被积函数虽在前半段为正、总积分仍可能为负的原因。
\end{exercise}
\begin{exercise}{627}{分部积分与端点数据}
设 $f\in C^2[0,2]$，$f(0)=1,f(2)=3,f'(2)=5$。求 $\int_0^2xf''(x)\dd x$，说明为什么无需知道 $f'(0)$。若再知 $f'(0)=a$，求 $\int_0^2(2-x)f''(x)\dd x$。
\end{exercise}
\begin{exercise}{628}{周期积分的常值与符号}
设 $F(x)=\int_x^{x+2\pi}e^{\sin t}\sin t\dd t$。证明 $F$ 是常函数，并确定其符号，不能只根据因子 $\sin t$ 在一周期积分为零判断。用正、负半周配对给出严格证明。
\end{exercise}
\begin{exercise}{629}{已知原函数与数据化积分}
已知 $e^{x^2}$ 是 $f(x)$ 的一个原函数。求 $\int_1^e\frac{\ln x\,f'(\ln x)}x\dd x$；计算中先换元，再用分部积分把结果化为原函数和 $f$ 的端点值，说明所需可导性由何保证。
\end{exercise}
\begin{exercise}{630}{有限秩积分方程}
设连续函数 $f$ 满足 $f(x)=x+2\int_0^{\pi/2}f(t)\cos t\dd t$。求 $f$ 及 $\int_0^1f(x)\dd x$。若积分前的系数 $2$ 改为实数 $\lambda$，讨论连续解的存在性和唯一性。
\end{exercise}
\begin{exercise}{631}{可去端点与积分次序转换}
设 $f(x)=\int_0^x\frac{\sin t}{\pi-t}\dd t$，$0\le x\le\pi$，在 $t=\pi$ 处按连续延拓取值。求 $\int_0^\pi f(x)\dd x$；先证明延拓存在，再用分部积分或交换累次积分次序计算，写明边界项。
\end{exercise}
\begin{exercise}{632}{变限积分的分部计算}
设 $f(x)=\int_1^xe^{-y^2}\dd y$。求 $\int_0^1x^2f(x)\dd x$，说明虽然 $e^{-x^2}$ 的原函数不能用初等函数表示，本题仍可通过分部积分化为初等积分；由 $f$ 的符号核验结果。
\end{exercise}
\begin{exercise}{633}{含导数积分方程与积分计算}
设 $f\in C^1[0,\pi/2]$，满足 $f(x)=x-\int_0^xf'(t)\cos t\dd t$。求 $f$ 及 $\int_0^{\pi/2}f(x)\dd x$，并代回原等式验证端点条件与函数表达式。
\end{exercise}
\begin{exercise}{634}{对称换元与参数积分}
对 $a>0$，计算
\[
I(a)=\int_0^{\pi/2}\frac{\sin x}{\sin x+a\cos x}\dd x.
\]
利用 $x\mapsto\frac\pi2-x$ 建立 $I(a)$ 与 $I(1/a)$ 的关系，并单独处理 $a=1$ 的情形。
\end{exercise}
\begin{exercise}{635}{分部积分与递推估计}
记 $I_n=\int_0^1x^ne^x\dd x$，$n=0,1,2,\ldots$。建立 $I_n$ 的递推关系，求 $I_3$，并证明 $\lim_{n\to\infty}(n+1)I_n=e$。
\end{exercise}
\subsection{反常积分}
\begin{exercise}{636}{无穷积分分类与判敛}
取 $a>1$，分别判断 $\int_a^{+\infty}x^{-1/2}\dd x$、$\int_a^{+\infty}\frac{\ln x}{x}\dd x$、$\int_a^{+\infty}\frac{\dd x}{x\ln x}$、$\int_a^{+\infty}xe^{-x}\dd x$ 的敛散性，给出理由；对收敛者求值。
\end{exercise}
\begin{exercise}{637}{对数换元与参数临界值}
计算 $\int_2^{+\infty}\frac{\dd x}{x\ln^2x}$。对 $a>1$、实参数 $p$，讨论 $\int_a^{+\infty}\frac{\dd x}{x\ln^px}$ 的敛散性，收敛时求值。
\end{exercise}
\begin{exercise}{638}{全轴反常积分与平移配方}
计算 $\int_{-\infty}^{+\infty}\frac{\dd x}{x^2+4x+5}$，写明全轴积分的定义。再求使 $\int_{-\infty}^{+\infty}\frac{\dd x}{x^2+4x+c}$ 收敛的实参数 $c$ 的范围，并在收敛时求值。
\end{exercise}
\begin{exercise}{639}{拆项积分与无穷端点抵消}
计算 $\int_1^{+\infty}\frac{\dd x}{x^2+x}$。说明为何部分分式分解后不能将积分写为两个发散积分之差；再求下限改为 $a>0$ 时的值及 $a\to0^+$ 时的极限。
\end{exercise}
\begin{exercise}{640}{反常根式积分与三角代换}
计算 $\int_3^{+\infty}\frac{\dd x}{(x-1)^4\sqrt{x^2-2x}}$。先用 $u=x-1$ 化简，再选择三角代换；求下限改为 $a>2$ 时的值，并判断 $a\to2^+$ 时的极限是否有限。
\end{exercise}
\begin{exercise}{641}{全轴换元与收敛证明}
计算 $\int_{-\infty}^{+\infty}\frac{e^x}{1+e^{2x}}\dd x$，分别说明两个无穷端点的收敛性。将指数中的 $x$ 换成 $ax$（$a\ne0$），求积分并讨论 $a$ 的符号。
\end{exercise}
\begin{exercise}{642}{分段积分与参数阈值}
计算 $\int_0^{+\infty}\min\{e^{-x},1/2\}\dd x$。将 $1/2$ 改为 $c>0$，分情况求积分关于 $c$ 的表达式，并核验其在 $c=1$ 连续。
\end{exercise}
\begin{exercise}{643}{反常分部积分与边界项}
计算 $\int_1^{+\infty}\frac{\ln x}{(1+x)^2}\dd x$。先证明积分收敛并处理分部积分中的无穷端边界项，再求下限改为 $a\ge1$ 时的积分。
\end{exercise}
\begin{exercise}{644}{三角换元与全轴推广}
计算 $\int_0^{+\infty}\frac{\dd x}{(1+x^2)^2}$。进一步求 $\int_{-\infty}^{+\infty}\frac{\dd x}{(a^2+x^2)^2}$（$a>0$），说明尺度代换中的系数。
\end{exercise}
\begin{exercise}{645}{反常三角换元}
计算 $\int_1^{+\infty}\frac{\dd x}{x\sqrt{2x^2-1}}$。再讨论下限改为 $a\ge1/\sqrt2$ 时的收敛性并求值，特别说明临界下限处的瑕点是否可积。
\end{exercise}
\begin{exercise}{646}{反常分部积分与对数积分}
计算 $\int_0^{+\infty}\frac{xe^{-x}}{(1+e^{-x})^2}\dd x$。先用比较法证明收敛，再选择使无穷端边界项为零的分部方式；最后计算将 $x$ 同时替换为 $ax$（$a>0$）后的积分。
\end{exercise}
\begin{exercise}{647}{奇函数与反常积分定义}
判断 $\int_{-\infty}^{+\infty}\sin2x\,e^{|x|}\dd x$ 的敛散性，并与对称截断积分 $\lim_{R\to+\infty}\int_{-R}^R\sin2x\,e^{|x|}\dd x$ 比较。证明两者不能仅由奇偶性等同；若 $e^{|x|}$ 改为 $e^{-|x|}$，判断是否绝对收敛并求积分值。
\end{exercise}
\begin{exercise}{648}{比较判别法与发散证明}
判断 $\int_2^{+\infty}\frac{\dd x}{\ln x}$ 的敛散性，给出明确的比较函数与成立区间。再讨论 $\int_2^{+\infty}\frac{\dd x}{x(\ln x)^p}$ 随实数 $p$ 的敛散性，说明增加 $x$ 因子后的变化。
\end{exercise}
\begin{exercise}{649}{高斯型积分的收敛性}
不用高斯积分的准确值，证明 $\int_0^{+\infty}e^{-x^2}\dd x$ 收敛。对 $a\in\mathbb R$ 讨论 $\int_0^{+\infty}e^{-ax^2}\dd x$ 的敛散性，并在 $a>0$ 时给出它与 $a=1$ 时积分的关系。
\end{exercise}
\begin{exercise}{650}{指数型极限与反常积分方程}
求实数 $c$，使 $\lim_{x\to+\infty}\left(\frac{x+c}{x-c}\right)^x=\int_{-\infty}^cte^{2t}\dd t$。分别计算两边，说明极限中幂底对充分大的 $x$ 为正，并证明参数解唯一。
\end{exercise}
\begin{exercise}{651}{瑕点识别与可去间断点}
分别考察 $\int_0^1\frac{\sin x}{x}\dd x$、$\int_0^1x\ln x\dd x$、$\int_{-1}^1\frac{\dd x}{\sqrt{x^2}}$、$\int_0^1x^{-2}\dd x$。指出未定义点，区分可作连续延拓的情形与无界瑕点，判断敛散性，并计算收敛积分中 $\int_0^1x\ln x\dd x$ 的值。
\end{exercise}
\begin{exercise}{652}{对数端点与连续延拓}
判断 $x\ln^2x$ 在 $x=0$ 是否无界；通过连续延拓计算 $\int_0^1x\ln^2x\dd x$。推广讨论 $\int_0^1x^p\ln^2x\dd x$ 随实参数 $p$ 的敛散性，收敛时求值。
\end{exercise}
\begin{exercise}{653}{端点瑕积分与参数条件}
对 $a>0$，计算 $\int_0^a\frac{\dd x}{\sqrt{a^2-x^2}}$，用极限说明上端点的处理。再讨论 $\int_0^a(a^2-x^2)^{-p}\dd x$ 随实参数 $p$ 的敛散性。
\end{exercise}
\begin{exercise}{654}{多瑕点与对数换元}
判断 $\int_0^1\frac{\dd x}{x(\ln x)^2}$ 的敛散性，分别考察 $0$ 和 $1$；再判断 $\int_0^{+\infty}\frac{\dd x}{x(\ln x)^2}$ 是否存在，指出全部需要拆分的瑕点，不能只检查无穷远。
\end{exercise}
\begin{exercise}{655}{瑕积分的幂次判别}
设 $a<b$，判断 $\int_a^b\frac{\dd x}{(x-a)^2}$ 的敛散性。推广讨论 $\int_a^b\frac{\dd x}{(x-a)^p}$（$p\in\mathbb R$），收敛时求值，并明确临界指数。
\end{exercise}
\begin{exercise}{656}{倒数换元与对数反常积分}
证明 $\int_0^{+\infty}\frac{\ln x}{1+x^2}\dd x$ 绝对收敛，再用 $x=1/t$ 求其值。说明在收敛性尚未建立前，不能直接依换元得到 $I=-I$ 就下结论。
\end{exercise}
\begin{exercise}{657}{分段反常积分与双重条件}
设 $f(x)=(x-1)^{1-\alpha}$（$1<x<e$），$f(x)=1/[x\ln^{\alpha+1}x]$（$x\ge e$）。求使 $\int_1^{+\infty}f(x)\dd x$ 收敛的全部实数 $\alpha$，并在收敛范围内求积分值。说明两个端点分别给出什么参数条件。
\end{exercise}
\begin{exercise}{658}{反常积分运算的有效条件}
分别判断 $\int_{-\infty}^{+\infty}\frac{x}{1-x^2}\dd x$、$\int_{-\infty}^{+\infty}\sin x\dd x$、$\int_{-1}^1x^{-2}\dd x$ 是否存在，指出用奇偶性或直接代原函数端点的失效条件。再计算 $\int_0^{+\infty}x^{3/2}e^{-x}\dd x$，可使用 $\Gamma(s+1)=s\Gamma(s)$、$\Gamma(1/2)=\sqrt\pi$，并说明其收敛性。
\end{exercise}
\begin{exercise}{659}{双端瑕点与参数判敛}
对实数 $p,q$，讨论反常积分
\[
\int_0^1\frac{\dd x}{x^p(1-x)^q}
\]
的敛散性，并在 $p=q=\frac12$ 时计算其值。说明为何必须分别考察两个端点。
\end{exercise}
\begin{exercise}{660}{对数因子与临界指数}
讨论 $\int_e^{+\infty}\frac{\dd x}{x(\ln x)^p}$ 与 $\int_0^{1/e}\frac{\dd x}{x|\ln x|^p}$ 随实参数 $p$ 的敛散性；收敛时分别求值，并用换元解释二者的联系。
\end{exercise}
